\item Para minimizar las fugas de neutrones en el núcleo de un reactor nuclear, la relación entre el área superficial de confinamiento y el volumen del reactor debe ser mínima. Para un volumen fijo dado $V$, calcule la relación de área/volumen para las siguientes geometrías:
\begin{enumerate}
    \item una esfera de radio $R$,
    \item un cubo de lado $L$, y
    \item un paralelepípedo de dimensiones $a \times a \times 2a$.
    \item ¿Cuál de estas formas maximiza y cuál minimiza las fugas de neutrones? Justifique cuantitativamente.
\end{enumerate}